\documentclass[12pt]{article}
\usepackage{amsmath, amssymb, geometry}
\geometry{margin=1in}
\setlength{\parindent}{0pt}
\setlength{\parskip}{0.6em}

\begin{document}

\begin{center}
\Large \textbf{Solutions -- Quiz 3\\ \vspace{0.6em}ECON 300 -- Labour Economics}\\

\vspace{0.6em}
\normalsize \textbf{March 04, 2026}
\end{center}

\vspace{0.5em}

\textbf{Part A: Multiple Choice (8 marks; 2 marks each)}\\
Choose the \textbf{one best answer} and provide a brief justification.

\vspace{0.6em}

\textbf{1.} \textbf{Correct answer: B}\\
\textbf{Explanation:} A monopsonist faces an upward-sloping labour supply curve. To hire one more worker, it typically must raise the wage, which increases the wage paid not only to the marginal worker but also to existing workers (in the standard single-wage monopsony). Therefore, the \textbf{marginal cost of labour exceeds the wage}, meaning the cost of recruiting an additional worker involves a \textbf{larger increase in total labour costs} than under perfect competition (where the firm is a wage taker and can hire at a constant wage).

\vspace{0.6em}

\textbf{2.} \textbf{Correct answer: A}\\
\textbf{Explanation:} The false statement is that the monopsonist is a \textbf{wage taker}. In monopsony, the firm has labour market power and faces an upward-sloping supply of labour; it can influence the wage by changing employment. The other statements are true in the standard model: the monopsonist has market power, is profit-maximizing, and may practice wage discrimination. Also, ``technically speaking'' the monopsonist does not take a wage as given and instead chooses employment given the supply curve; its labour demand is not represented as a standard competitive demand curve in the same sense.

\vspace{0.6em}

\textbf{3.} \textbf{Correct answer: E}\\
\textbf{Explanation:} In monopsony, the marginal cost of labour (MC) lies \textbf{above} the average cost of labour (AC), where AC is the wage on the labour supply curve. This occurs because raising employment requires raising wages, increasing the wage bill on inframarginal workers. Therefore, it is false that MC is lower than AC. The other statements are standard monopsony results: wages and employment are lower than competitive levels, worker rents are reduced (transferred to the employer), and the total wage bill may be lower relative to competition.

\vspace{0.6em}

\textbf{4.} \textbf{Correct answer: B}\\
\textbf{Explanation:} A firm in a perfectly competitive labour market is a \textbf{wage taker} and faces a \textbf{perfectly elastic (infinitely elastic) labour supply} at the market wage. It does \emph{not} need to operate in a perfectly competitive output market, and worker unionization is not required for perfect competition in labour markets.

\newpage

\textbf{Part B: Problem (12 marks)}

\textbf{Given:}
\[
N^S = 25W,
\qquad
N^D = 500 - 12.5W,
\]
where $W$ is wage and $N$ is employment.

\vspace{0.4em}

\textbf{(a) Competitive equilibrium (4 marks)}\\
Equilibrium occurs where $N^S=N^D$:
\[
25W = 500 - 12.5W.
\]
Combine like terms:
\[
25W + 12.5W = 500
\Rightarrow 37.5W = 500
\Rightarrow W^* = \frac{500}{37.5}.
\]
Compute:
\[
W^* = \frac{500}{37.5} = \frac{5000}{375} = \frac{40}{3}\approx 13.33.
\]
Now compute employment:
\[
N^* = 25W^* = 25\left(\frac{40}{3}\right)=\frac{1000}{3}\approx 333.33.
\]

\vspace{0.6em}

\textbf{Answer:} 
\[
W^*=\frac{40}{3}\approx 13.33,
\qquad
N^*=\frac{1000}{3}\approx 333.33.
\]

\vspace{0.8em}

\textbf{(b) Worker surplus and firm surplus (4 marks)}\\
We are instructed to use \textbf{triangle areas} (no calculus).

\textbf{Step 1: Inverse supply and inverse demand.}\\
Supply:
\[
N=25W \Rightarrow W_S(N)=\frac{N}{25}.
\]
This passes through the origin, so the supply intercept is $W=0$.

Demand:
\[
N=500-12.5W
\Rightarrow 12.5W = 500-N
\Rightarrow W_D(N)=\frac{500-N}{12.5}.
\]
The demand intercept (when $N=0$) is:
\[
W_D(0)=\frac{500}{12.5}=40.
\]

\textbf{Step 2: Worker surplus (WS).}\\
WS is the area between the wage line $W^*$ and the supply curve up to $N^*$.
Since the supply intercept is 0, WS is a triangle with:
\[
\text{base}=N^*=\frac{1000}{3}, 
\qquad
\text{height}=W^*=\frac{40}{3}.
\]
Thus:
\[
WS=\frac{1}{2}\cdot \text{base}\cdot \text{height}
= \frac{1}{2}\left(\frac{1000}{3}\right)\left(\frac{40}{3}\right)
= \frac{20000}{9}
\approx 2222.22.
\]

\textbf{Step 3: Firm surplus (FS).}\\
FS is the area between the demand curve and the wage line $W^*$ up to $N^*$.
This is also a triangle with:
\[
\text{base}=N^*=\frac{1000}{3},
\qquad
\text{height}=\left(40-W^*\right)=40-\frac{40}{3}=\frac{80}{3}.
\]
Thus:
\[
FS=\frac{1}{2}\left(\frac{1000}{3}\right)\left(\frac{80}{3}\right)
= \frac{40000}{9}
\approx 4444.44.
\]

\vspace{0.6em}

\textbf{Answer:}
\[
WS=\frac{20000}{9}\approx 2222.22,
\qquad
FS=\frac{40000}{9}\approx 4444.44.
\]

\vspace{0.8em}

\textbf{(c) Demand increases to $N^{D'}=600-12.5W$ (4 marks)}\\
New equilibrium is where $N^S=N^{D'}$:
\[
25W = 600 - 12.5W.
\]
Combine like terms:
\[
25W + 12.5W = 600
\Rightarrow 37.5W=600
\Rightarrow W_1 = \frac{600}{37.5}.
\]
Compute:
\[
W_1 = \frac{600}{37.5} = \frac{6000}{375} = 16.
\]
Now compute employment:
\[
N_1 = 25W_1 = 25(16)=400.
\]

\textbf{Explanation (words):} The increase in labour demand shifts the demand curve outward. At the original wage, firms wish to hire more labour than workers supply, creating upward pressure on wages. As wages rise, more workers supply labour and firms move along the new demand curve, resulting in a higher equilibrium wage and higher equilibrium employment.

\vspace{0.6em}

\textbf{Answer:}
\[
W_1 = 16,
\qquad
N_1 = 400.
\]

\end{document}